
\chapter{2016年北京卷（理）}

\section{单选题}

\begin{problem}
已知集合 \(A=\{x\mid |x|<2\}\)，\(B=\{-1,0,1,2,3\}\)，则 \(A\cap B=\)（\quad）

\choices
    {\(\{0,1\}\)}
    {\(\{0,1,2\}\)}
    {\(\{-1,0,1\}\)}
    {\(\{-1,0,1,2\}\)}
\end{problem}

\begin{answer}
C
\end{answer}

\begin{solution}
由 \(|x|<2\) 得 \(-2<x<2\).集合 \(B\) 中只有 \(-1\)，\(0\)，\(1\) 满足这个不等式，\(2\) 和 \(3\) 均不满足.因此
\(A\cap B=\{-1,0,1\}\)，故选 C.
\end{solution}

\begin{problem}
若 \(x\)，\(y\) 满足 \(\begin{cases}
2x-y\le 0,\\
x+y\le 3,\\
x\ge 0,
\end{cases}\)
则 \(2x+y\) 的最大值为（\quad）

\choices
    {\(0\)}
    {\(3\)}
    {\(4\)}
    {\(5\)}
\end{problem}

\begin{answer}
C
\end{answer}

\begin{solution}
由 \(2x-y\le 0\) 得 \(y\ge 2x\).在 \(x\ge0\)，\(y\ge2x\)，\(x+y\le3\) 所确定的可行域中，边界交点为 \((0,0)\)，\((0,3)\)，\((1,2)\).目标函数 \(2x+y\) 在这三个点处的值分别为 \(0\)，\(3\)，\(4\)，所以最大值为 \(4\)，故选 C.
\end{solution}

\begin{problem}
执行如图所示的程序框图，若输入的 \(a\) 值为 \(1\)，则输出的 \(k\) 值为（\quad）

\texfigure[max width=6.0cm]{img_repaint/2016/beijing_science/q3_fig1.tex}

\choices
    {\(1\)}
    {\(2\)}
    {\(3\)}
    {\(4\)}
\end{problem}

\begin{answer}
B
\end{answer}

\begin{solution}
输入 \(a=1\) 后，初始值为 \(k=0\)，\(b=1\).第一次循环得到 \(a=-1/(1+1)=-1/2\)，因 \(a\ne b\)，令 \(k=1\)；第二次循环得到 \(a=-1/(1-1/2)=-2\)，因 \(a\ne b\)，令 \(k=2\)；第三次循环得到 \(a=-1/(1-2)=1=b\)，此时不再增加 \(k\)，直接输出 \(k=2\)，故选 B.
\end{solution}

\begin{problem}
设 \(\bs a\)，\(\bs b\) 是向量，则“\(|\bs a|=|\bs b|\)”是“\(|\bs a+\bs b|=|\bs a-\bs b|\)”的（\quad）

\choices
    {充分而不必要条件}
    {必要而不充分条件}
    {充分必要条件}
    {既不充分也不必要条件}
\end{problem}

\begin{answer}
D
\end{answer}

\begin{solution}
将两个等式分别平方，得
\(|\bs a+\bs b|^2=|\bs a|^2+|\bs b|^2+2\bs a\cdot\bs b\)，
\(|\bs a-\bs b|^2=|\bs a|^2+|\bs b|^2-2\bs a\cdot\bs b\).因此
\(|\bs a+\bs b|=|\bs a-\bs b|\) 等价于 \(\bs a\cdot\bs b=0\).

取 \(\bs a=\bs b=(1,0)\)，则 \(|\bs a|=|\bs b|\)，但 \(|\bs a+\bs b|=2\ne0=|\bs a-\bs b|\)，所以前者不是后者的充分条件.取 \(\bs a=(1,0)\)，\(\bs b=(0,2)\)，则 \(\bs a\cdot\bs b=0\)，后者成立，但 \(|\bs a|=1\ne2=|\bs b|\)，所以前者也不是后者的必要条件.故选 D.
\end{solution}

\begin{problem}
已知 \(x\)，\(y\in\R\)，且 \(x>y>0\)，则（\quad）

\choices
    {\(\dfrac1x-\dfrac1y>0\)}
    {\(\sin x-\sin y>0\)}
    {\(\left(\dfrac12\right)^x-\left(\dfrac12\right)^y<0\)}
    {\(\ln x+\ln y>0\)}
\end{problem}

\begin{answer}
C
\end{answer}

\begin{solution}
因为 \(x>y>0\)，所以 \(xy>0\)，从而
\(\dfrac1x-\dfrac1y=\dfrac{y-x}{xy}<0\)，选项 A 错.

正弦函数在 \((0,+\infty)\) 上不是单调函数，例如取 \(x=\pi\)，\(y=\dfrac\pi2\)，有 \(\sin x-\sin y=-1<0\)，故选项 B 错.函数 \(y=\left(\dfrac12\right)^x\) 在 \(\R\) 上单调递减，所以 \(x>y\) 时 \(\left(\dfrac12\right)^x<\left(\dfrac12\right)^y\)，选项 C 正确.

又 \(\ln x+\ln y=\ln(xy)\)，而 \(xy\) 不一定大于 \(1\)，例如取 \(x=\dfrac12\)，\(y=\dfrac14\) 时该和小于 \(0\)，故选项 D 错.因此选 C.
\end{solution}

\begin{problem}
某三棱锥的三视图如图所示，则该三棱锥的体积为（\quad）

\texfigure[max width=0.48\linewidth]{img_repaint/2016/beijing_science/q6_fig1.tex}

\choices
    {\(\dfrac16\)}
    {\(\dfrac13\)}
    {\(\dfrac12\)}
    {\(1\)}
\end{problem}

\begin{answer}
A
\end{answer}

\begin{solution}
由俯视图可知，该三棱锥的底面是两条直角边长均为 \(1\) 的直角三角形；由正视图和侧视图可知棱锥的高为 \(1\).因此底面积
\(S=\dfrac12\times1\times1=\dfrac12\)，体积
\(V=\dfrac13Sh=\dfrac13\times\dfrac12\times1=\dfrac16\)，故选 A.
\end{solution}

\begin{problem}
将函数 \(y=\sin\left(2x-\dfrac{\pi}{3}\right)\) 图象上的点 \(P\left(\dfrac{\pi}{4},t\right)\) 向左平移 \(s\)（\(s>0\)）个单位长度得到点 \(P'\)，若 \(P'\) 位于函数 \(y=\sin2x\) 的图象上，则（\quad）

\choices
    {\(t=\dfrac12\)，\(s\) 的最小值为 \(\dfrac{\pi}{6}\)}
    {\(t=\dfrac{\sqrt3}{2}\)，\(s\) 的最小值为 \(\dfrac{\pi}{6}\)}
    {\(t=\dfrac12\)，\(s\) 的最小值为 \(\dfrac{\pi}{3}\)}
    {\(t=\dfrac{\sqrt3}{2}\)，\(s\) 的最小值为 \(\dfrac{\pi}{3}\)}
\end{problem}

\begin{answer}
A
\end{answer}

\begin{solution}
点 \(P\) 在函数 \(y=\sin\left(2x-\dfrac\pi3\right)\) 的图象上，所以
\(t=\sin\left(\dfrac\pi2-\dfrac\pi3\right)=\sin\dfrac\pi6=\dfrac12\).

向左平移后，\(P'\) 的坐标为 \(\left(\dfrac\pi4-s,\dfrac12\right)\).由于 \(P'\) 在 \(y=\sin2x\) 的图象上，有
\(\sin\left(\dfrac\pi2-2s\right)=\dfrac12\).因此
\(\dfrac\pi2-2s=\dfrac\pi6+2k\pi\) 或 \(\dfrac{5\pi}6+2k\pi\)，其中 \(k\in\Z\)，即
\(s=\dfrac\pi6-k\pi\) 或 \(s=-\dfrac\pi6-k\pi\).在 \(s>0\) 的条件下，第一类的最小值为 \(\dfrac\pi6\)，第二类的最小正值为 \(\dfrac{5\pi}6\).所以 \(t=\dfrac12\)，且 \(s\) 的最小值为 \(\dfrac\pi6\)，故选 A.
\end{solution}

\begin{problem}
袋中装有偶数个球，其中红球，黑球各占一半. 甲，乙，丙是三个空盒. 每次从袋中任意取出两个球，将其中一个球放入甲盒，如果这个球是红球，就将另一个放入乙盒，否则就放入丙盒. 重复上述过程，直到袋中所有球都被放入盒中，则（\quad）

\choices
    {乙盒中黑球不多于丙盒中黑球}
    {乙盒中红球与丙盒中黑球一样多}
    {乙盒中红球不多于丙盒中红球}
    {乙盒中黑球与丙盒中红球一样多}
\end{problem}

\begin{answer}
B
\end{answer}

\begin{solution}
设袋中红球、黑球各有 \(m\) 个，于是共进行 \(m\) 次取球.记四类取球结果的次数分别为：第一球红、第二球红的次数 \(r_{RR}\)；第一球黑、第二球黑的次数 \(r_{BB}\)；第一球红、第二球黑的次数 \(r_{RB}\)；第一球黑、第二球红的次数 \(r_{BR}\).

按颜色统计所有球，红球数和黑球数相等，故
\(2r_{RR}+r_{RB}+r_{BR}=2r_{BB}+r_{RB}+r_{BR}\)，从而 \(r_{RR}=r_{BB}\).第一球为红时第二球放入乙盒，所以乙盒中的红球数为 \(r_{RR}\)；第一球为黑时第二球放入丙盒，所以丙盒中的黑球数为 \(r_{BB}\).二者相等，故选 B.
\end{solution}

\section{填空题}

\begin{problem}
设 \(a\in\R\)，若复数 \((1+\i)(a+\i)\) 在复平面内对应的点位于实轴上，则 \(a=\fillinblank{}\).
\end{problem}

\begin{answer}
\(-1\)
\end{answer}

\begin{solution}
\((1+\i)(a+\i)=a+\i+a\i+\i^2=(a-1)+(a+1)\i\).对应点位于实轴上，当且仅当虚部为零，即 \(a+1=0\)，所以 \(a=-1\).
\end{solution}

\begin{problem}
在 \((1-2x)^6\) 的展开式中，\(x^2\) 的系数为\fillinblank{}（用数字作答）.
\end{problem}

\begin{answer}
60
\end{answer}

\begin{solution}
由二项式定理，\((1-2x)^6\) 中含 \(x^2\) 的项来自取两次 \(-2x\)，其系数为
\(\mathrm C_6^2(-2)^2=15\times4=60\).
\end{solution}

\begin{problem}
在极坐标系中，直线 \(\rho\cos\theta-\sqrt{3}\rho\sin\theta-1=0\) 与圆 \(\rho=2\cos\theta\) 交于 \(A\)，\(B\) 两点，则 \(|AB|=\fillinblank{}\).
\end{problem}

\begin{answer}
2
\end{answer}

\begin{solution}
令 \(x=\rho\cos\theta\)，\(y=\rho\sin\theta\).直线方程化为 \(x-\sqrt3y=1\)，圆方程 \(\rho=2\cos\theta\) 化为 \(x^2+y^2=2x\)，即圆心为 \((1,0)\)，半径为 \(1\).

圆心 \((1,0)\) 满足直线方程，因此该直线经过圆心，弦 \(AB\) 是直径.所以 \(|AB|=2\).
\end{solution}

\begin{problem}
已知 \(\{a_n\}\) 为等差数列，\(S_n\) 为其前 \(n\) 项和. 若 \(a_1=6\)，\(a_3+a_5=0\)，则 \(S_6=\fillinblank{}\).
\end{problem}

\begin{answer}
6
\end{answer}

\begin{solution}
设等差数列的公差为 \(d\)，则 \(a_3=6+2d\)，\(a_5=6+4d\).由 \(a_3+a_5=0\) 得
\(12+6d=0\)，所以 \(d=-2\).因此
\(S_6=\dfrac{6}{2}\bigl(2a_1+5d\bigr)=3\bigl(12-10\bigr)=6\).
\end{solution}

\begin{problem}
双曲线 \(\dfrac{x^2}{a^2}-\dfrac{y^2}{b^2}=1\)（\(a>0\)，\(b>0\)）的渐近线为正方形 \(OABC\) 的边 \(OA\)，\(OC\) 所在的直线，点 \(B\) 为该双曲线的焦点. 若正方形 \(OABC\) 的边长为 \(2\)，则 \(a=\fillinblank{}\).
\end{problem}

\begin{answer}
2
\end{answer}

\begin{solution}
双曲线的渐近线为 \(y=\pm\dfrac ba x\).它们分别是正方形从顶点 \(O\) 引出的两条边所在的直线，而正方形的相邻边垂直，所以两条渐近线垂直，从而 \(\dfrac ba=1\)，即 \(a=b\).

正方形边长为 \(2\)，取两条渐近线方向为 \((1,1)\) 和 \((1,-1)\)，则两条边向量可写为 \((\sqrt2,\sqrt2)\) 和 \((\sqrt2,-\sqrt2)\)，故对角顶点 \(B\) 为 \((2\sqrt2,0)\)，有焦距 \(c=OB=2\sqrt2\).又 \(c^2=a^2+b^2=2a^2\)，所以 \(2a^2=8\)，且 \(a>0\)，得 \(a=2\).
\end{solution}

\begin{problem}
设函数
\(f(x)=\begin{cases}
x^3-3x, & x\le a,\\
-2x, & x>a.
\end{cases}\)

\begin{enumerate}
\item 若 \(a=0\)，则 \(f(x)\) 的最大值为\fillinblank{}；
\item 若 \(f(x)\) 无最大值，则实数 \(a\) 的取值范围是\fillinblank{}.
\end{enumerate}
\end{problem}

\begin{answer}
\(2\)，\((-\infty,-1)\)
\end{answer}

\begin{solution}
当 \(a=0\) 时，\(x\le0\) 上 \(f(x)=x^3-3x\)，其导数为 \(f'(x)=3x^2-3\).因此在 \((-\infty,-1)\) 上递增，在 \((-1,0)\) 上递减，且 \(f(-1)=2\)，而 \(x>0\) 时 \(f(x)=-2x<0\).所以函数最大值为 \(2\).

一般地，记 \(g(x)=x^3-3x\)，考察左段 \((-\infty,a]\) 的最大值.若 \(a\le-1\)，则 \(g\) 在该区间上递增，左段最大值为 \(g(a)\)；若 \(a>-1\)，左段至少包含点 \(-1\)，且其最大值存在，记为 \(M(a)\).右段 \(x>a\) 上 \(f(x)=-2x\) 严格递减，其上确界为 \(-2a\)，但在右段不能取到.

当 \(a<-1\) 时，\(-2a-g(a)=-2a-(a^3-3a)=a-a^3=-a(a-1)(a+1)>0\)，所以右段的上确界严格大于左段最大值，函数没有最大值.\(a=-1\) 时左段在 \(x=-1\) 处取到值 \(2\)，而右段值均小于 \(2\)，故仍有最大值.\(a>-1\) 时左段最大值至少为 \(2\)，而 \(-2a<2\)，故函数有最大值.因此无最大值时 \(a\in(-\infty,-1)\).
\end{solution}

\section{解答题}

\begin{problem}
在 \(\triangle ABC\) 中，\(a^2+c^2=b^2+\sqrt{2}ac\).

\begin{enumerate}
\item 求 \(\angle B\) 的大小；
\item 求 \(\sqrt{2}\cos A+\cos C\) 的最大值.
\end{enumerate}
\end{problem}

\begin{solution}
\begin{enumerate}
\item 由余弦定理，\(b^2=a^2+c^2-2ac\cos B\).与题设等式比较，得 \(2ac\cos B=\sqrt2ac\).由于三角形边长 \(a\)，\(c>0\)，所以 \(\cos B=\dfrac{\sqrt2}{2}\).又 \(0<B<\pi\)，故 \(\angle B=\dfrac\pi4\).
\item 由三角形内角和，\(A+C=\pi-B=\dfrac{3\pi}{4}\).令 \(C=\dfrac{3\pi}{4}-A\)，其中 \(0<A<\dfrac{3\pi}{4}\)，则
\(\sqrt2\cos A+\cos C=\sqrt2\cos A+\cos\left(\dfrac{3\pi}{4}-A\right)=\dfrac{\sqrt2}{2}(\cos A+\sin A)\).
因为 \(\cos A+\sin A\le\sqrt2\)，且当 \(A=\dfrac\pi4\)，\(C=\dfrac\pi2\) 时取等号，所以所求最大值为 \(1\).
\end{enumerate}
\end{solution}

\begin{problem}
\(A\)，\(B\)，\(C\) 三个班共有 \(100\) 名学生，为调查他们的体育锻炼情况，通过分层抽样获得了部分学生一周的锻炼时间，数据如表（单位：小时）：

\begin{center}
\begin{tabular}{|c|cccccccc|}
\hline
\(A\)班 & \(6\) & \(6.5\) & \(7\) & \(7.5\) & \(8\) & & & \\
\hline
\(B\)班 & \(6\) & \(7\) & \(8\) & \(9\) & \(10\) & \(11\) & \(12\) & \\
\hline
\(C\)班 & \(3\) & \(4.5\) & \(6\) & \(7.5\) & \(9\) & \(10.5\) & \(12\) & \(13.5\) \\
\hline
\end{tabular}
\end{center}

\begin{enumerate}
\item 试估计 \(C\) 班的学生人数；
\item 从 A 班和 C 班抽出的学生中，各随机选取一个人，A 班选出的人记为甲，C 班选出的人记为乙. 假设所有学生的锻炼时间相对独立，求该周甲的锻炼时间比乙的锻炼时间长的概率；
\item 再从 A，B，C 三班中各随机抽取一名学生，他们该周锻炼时间分别是 \(7\)，\(9\)，\(8.25\)（单位：小时），这 \(3\) 个新数据与表格中的数据构成的新样本的平均数记为 \(\mu_1\)，表格中数据的平均数记为 \(\mu_0\)，试判断 \(\mu_0\) 和 \(\mu_1\) 的大小（结论不要求证明）.
\end{enumerate}
\end{problem}

\begin{solution}
\begin{enumerate}
\item 表中从三个班抽出的学生共有 \(5+7+8=20\) 人，其中来自 \(C\) 班的有 \(8\) 人.分层抽样的抽样比为 \(20/100=1/5\)，所以估计 \(C\) 班学生人数为 \(8\div\dfrac15=40\).
\item 甲、乙分别从两个样本中等可能独立选取，因此共有 \(5\times8=40\) 个等可能的有序结果.甲的时间分别为 \(6\)，\(6.5\)，\(7\)，\(7.5\)，\(8\) 时，比它大的乙的时间个数分别为 \(2\)，\(3\)，\(3\)，\(3\)，\(4\)，故满足甲比乙长的结果数为 \(2+3+3+3+4=15\).所求概率为 \(\dfrac{15}{40}=\dfrac38\).
\item 表中数据的总和为
\(6+6.5+7+7.5+8+6+7+8+9+10+11+12+3+4.5+6+7.5+9+10.5+12+13.5=164\)，所以 \(\mu_0=\dfrac{164}{20}=\dfrac{41}{5}\).加入的三个数据之和为 \(7+9+8.25=24.25\)，故 \(\mu_1=\dfrac{164+24.25}{23}=\dfrac{753}{92}\).比较得 \(41\times92=3772>3765=753\times5\)，所以 \(\mu_0>\mu_1\)，即 \(\mu_1<\mu_0\).
\end{enumerate}
\end{solution}

\begin{problem}
如图，在四棱锥 \(P-ABCD\) 中，平面 \(PAD\perp\) 平面 \(ABCD\)，\(PA\perp PD\)，\(PA=PD\)，\(AB\perp AD\)，\(AB=1\)，\(AD=2\)，\(AC=CD=\sqrt5\).

\begin{enumerate}
\item 求证：\(PD\perp\) 平面 \(PAB\)；
\item 求直线 \(PB\) 与平面 \(PCD\) 所成角的正弦值；
\item 在棱 \(PA\) 上是否存在点 \(M\)，使得 \(BM\parallel\) 平面 \(PCD\)？若存在，求 \(\dfrac{AM}{AP}\) 的值，若不存在，说明理由.
\end{enumerate}

\texfigure[max width=0.55\linewidth]{img_repaint/2016/beijing_science/q17_fig2.tex}
\end{problem}

\begin{solution}
\begin{enumerate}
\item 建立如图所示的空间直角坐标系，取 \(A\) 为原点，令 \(AB\) 为 \(x\) 轴，\(AD\) 为 \(y\) 轴，底面为 \(z=0\).由已知可取
\(A=(0,0,0)\)，\(B=(1,0,0)\)，\(D=(0,2,0)\).设 \(C=(x,y,0)\)，由 \(AC=CD=\sqrt5\) 得 \(x^2+y^2=5\)，\(x^2+(y-2)^2=5\)，两式相减得 \(y=1\)，再由 \(x^2+y^2=5\) 得 \(x=\pm2\).结合四边形 \(ABCD\) 的顶点顺序，取 \(C=(2,1,0)\).
平面 \(PAD\perp\) 底面且含有 \(AD\)，故可设 \(P=(0,u,v)\).由 \(PA=PD\) 得 \(u^2+v^2=(u-2)^2+v^2\)，所以 \(u=1\)；再由 \(PA\perp PD\) 得 \((-1,-v)\cdot(1,-v)=0\)，所以 \(v^2=1\).取 \(P\) 在底面上方，得 \(P=(0,1,1)\).

于是 \(\overrightarrow{PD}=(0,1,-1)\)，\(\overrightarrow{PA}=(0,-1,-1)\)，\(\overrightarrow{PB}=(1,-1,-1)\).有 \(\overrightarrow{PD}\cdot\overrightarrow{PA}=0\)，且 \(\overrightarrow{PD}\cdot\overrightarrow{PB}=0\).直线 \(PA\)，\(PB\) 在平面 \(PAB\) 内相交于 \(P\)，故由线面垂直的判定定理，\(PD\perp\) 平面 \(PAB\).
\item 取平面 \(PCD\) 的一个法向量
\(\bs n=\overrightarrow{PC}\times\overrightarrow{PD}=(2,0,-1)\times(0,1,-1)=(1,2,2)\).直线 \(PB\) 的方向向量为 \(\bs d=\overrightarrow{PB}=(1,-1,-1)\)，因此直线与平面所成角 \(\theta\) 满足
\(\sin\theta=\dfrac{|\bs d\cdot\bs n|}{|\bs d|\,|\bs n|}=\dfrac{|-3|}{\sqrt3\times3}=\dfrac{\sqrt3}{3}\).
\item 设 \(M\) 在棱 \(PA\) 上，令 \(\dfrac{AM}{AP}=t\)，则 \(0\le t\le1\)，且 \(M=(0,t,t)\).于是 \(\overrightarrow{BM}=(-1,t,t)\).直线 \(BM\) 平行于平面 \(PCD\) 当且仅当其方向向量垂直于法向量 \(\bs n\)，即
\((-1,t,t)\cdot(1,2,2)=-1+4t=0\).解得 \(t=\dfrac14\)，它属于 \([0,1]\).又 \(\overrightarrow{PB}\cdot\bs n=-3\ne0\)，故点 \(B\) 不在平面 \(PCD\) 上，所取直线不是包含于该平面的情形.因此存在这样的点 \(M\)，且 \(\dfrac{AM}{AP}=\dfrac14\).
\end{enumerate}
\end{solution}

\begin{problem}
设函数 \(f(x)=x\e^{a-x}+bx\)，曲线 \(y=f(x)\) 在点 \((2,f(2))\) 处的切线方程为 \(y=(\e-1)x+4\).

\begin{enumerate}
\item 求 \(a\)，\(b\) 的值；
\item 求 \(f(x)\) 的单调区间.
\end{enumerate}
\end{problem}

\begin{solution}
\begin{enumerate}
\item 切线在 \(x=2\) 处的纵坐标为 \(2(\e-1)+4=2\e+2\)，斜率为 \(\e-1\).又
\(f(2)=2\e^{a-2}+2b\)，\(f'(x)=(1-x)\e^{a-x}+b\)，所以
\(\e^{a-2}+b=\e+1\)，\(-\e^{a-2}+b=\e-1\).两式相加得 \(b=\e\)，代回得 \(\e^{a-2}=1\)，故 \(a=2\).
\item 代入 \(a=2\)，得 \(f'(x)=(1-x)\e^{2-x}+\e\).令 \(u=x-1\)，则
\(f'(x)=\e\bigl(1-u\e^{-u}\bigr)\).当 \(u\le0\) 时 \(u\e^{-u}\le0<1\)；当 \(u>0\) 时 \(\e^u>u\)，故 \(u\e^{-u}<1\).因此对任意 \(x\in\R\)，都有 \(f'(x)>0\)，所以 \(f(x)\) 在 \(\R\) 上单调递增，单调递减区间不存在.
\end{enumerate}
\end{solution}

\begin{problem}
已知椭圆 \(C:\dfrac{x^2}{a^2}+\dfrac{y^2}{b^2}=1\)（\(a>b>0\)）的离心率为 \(\dfrac{\sqrt3}{2}\)，\(A(a,0)\)，\(B(0,b)\)，\(O(0,0)\)，\(\triangle OAB\) 的面积为 \(1\).

\begin{enumerate}
\item 求椭圆 \(C\) 的方程；
\item 设 \(P\) 是椭圆 \(C\) 上一点，直线 \(PA\) 与 \(y\) 轴交于点 \(M\)，直线 \(PB\) 与 \(x\) 轴交于点 \(N\). 求证：\(|AN|\cdot|BM|\) 为定值.
\end{enumerate}

\texfigure[max width=0.55\linewidth]{img_repaint/2016/beijing_science/q17_fig1.tex}
\end{problem}

\begin{solution}
\begin{enumerate}
\item 设椭圆的半焦距为 \(c\).由离心率，\(c=\dfrac{\sqrt3}{2}a\)，从而
\(b^2=a^2-c^2=a^2-\dfrac34a^2=\dfrac14a^2\).又 \(\triangle OAB\) 的面积为 \(\dfrac12ab=1\)，且 \(a\)，\(b>0\)，所以 \(ab=2\).由 \(b=\dfrac a2\) 得 \(a=2\)，\(b=1\)，故椭圆方程为 \(\dfrac{x^2}{4}+y^2=1\).
\item 设 \(P=(x,y)\) 是椭圆上且不同于端点 \(A\)，\(B\) 的点.由直线 \(PA\) 与 \(y\) 轴的交点为 \(M\)，直线 \(PB\) 与 \(x\) 轴的交点为 \(N\)，可得
\(M=\left(0,\dfrac{2y}{2-x}\right)\)，\(N=\left(\dfrac{x}{1-y},0\right)\).因此
\(|BM|=\dfrac{|x+2y-2|}{|2-x|}\)，\(|AN|=\dfrac{|x+2y-2|}{|1-y|}\).

椭圆方程为 \(x^2+4y^2=4\)，且因 \(P\ne A\)，\(B\)，有 \(2-x>0\)，\(1-y>0\).由此
\((x+2y-2)^2-4(2-x)(1-y)=x^2+4y^2-4=0\)，所以
\(|AN|\cdot|BM|=\dfrac{(x+2y-2)^2}{(2-x)(1-y)}=4\).因此该乘积为定值 \(4\).
\end{enumerate}
\end{solution}

\begin{problem}
设数列 \(A\)：\(a_1\)，\(a_2\)，\(\ldots\)，\(a_N\)（\(N\ge2\)）. 如果对小于 \(n\)（\(2\le n\le N\)）的每个正整数 \(k\) 都有 \(a_k<a_n\)，则称 \(n\) 是数列 \(A\) 的一个“\(G\)时刻”，记 \(G(A)\) 是数列 \(A\) 的所有“\(G\)时刻”组成的集合.

\begin{enumerate}
\item 对数列 \(A\)：\(-2\)，\(2\)，\(-1\)，\(1\)，\(3\)，写出 \(G(A)\) 的所有元素；
\item 证明：若数列 \(A\) 中存在 \(a_n\) 使得 \(a_n>a_1\)，则 \(G(A)\ne\varnothing\)；
\item 证明：若数列 \(A\) 满足 \(a_n-a_{n-1}\le1\)（\(n=2\)，\(3\)，\(\ldots\)，\(N\)），则 \(G(A)\) 的元素个数不小于 \(a_N-a_1\).
\end{enumerate}
\end{problem}

\begin{solution}
\begin{enumerate}
\item 按定义逐项比较：\(a_2=2>a_1=-2\)，所以 \(2\) 是 \(G\) 时刻；\(a_3=-1\) 不大于前面的 \(2\)，\(a_4=1\) 不大于前面的 \(2\)，而 \(a_5=3\) 大于前四项.因此 \(G(A)=\{2,5\}\).
\item 在 \(a_1\)，\(a_2\)，\(\ldots\)，\(a_n\) 中取最大值，并令 \(j\) 为最大值第一次出现的下标.因为存在 \(a_n>a_1\)，这个最大值大于 \(a_1\)，故 \(j\ge2\)；按 \(j\) 的最小性，对每个 \(k<j\) 都有 \(a_k<a_j\)，所以 \(j\) 是 \(G\) 时刻.因此 \(G(A)\ne\varnothing\).
\item 令 \(M_j=\max\{a_1,a_2,\ldots,a_j\}\).当且仅当 \(j\) 是 \(G\) 时刻时，\(M_j>M_{j-1}\)，并且此时 \(M_j=a_j\).若 \(j\) 是 \(G\) 时刻，则 \(M_{j-1}\ge a_{j-1}\)，从而由题设
\(M_j-M_{j-1}=a_j-M_{j-1}\le a_j-a_{j-1}\le1\).

从 \(M_1=a_1\) 到 \(M_N\) 的增加只发生在 \(G\) 时刻，设 \(G(A)\) 有 \(m\) 个元素，则 \(M_N-a_1\le m\).又 \(a_N\le M_N\)，所以
\(a_N-a_1\le M_N-a_1\le m=|G(A)|\)，即 \(|G(A)|\ge a_N-a_1\).
\end{enumerate}
\end{solution}
